\documentclass[11pt]{article}
\usepackage[margin=1in]{geometry}
\usepackage[T1]{fontenc}
\usepackage{lmodern,booktabs}
\usepackage[hidelinks]{hyperref}
\usepackage{xurl}
\setlength{\parindent}{0pt}
\setlength{\parskip}{6pt}
\title{Hide Parked Vehicles in the Native Perception Scene}
\author{Pan Dynamics Collision Avoidance Research}
\date{October 2, 2026}
\begin{document}
\maketitle

\section{Configuration and behavior}
\path{config/carlaPerceptionScene.json} now defaults to
\texttt{hideParkedVehicles=true}. The local native capture driver reads this
configuration and disables environment objects labeled Car, Truck, Bus,
Motorcycle or Bicycle before spawning the experimental ego and target.
CARLA's \texttt{enable\_environment\_objects} removes these static objects
from rendering and collision queries, including ray-cast LiDAR.
The implementation was checked against the installed CARLA
\path{Util/ObjectRegister.cpp}. No image masking is applied.

The active local driver is
\path{simulation_output/carla_full_perception_20261002/captureNative.py}.
It accepts an optional \texttt{--scene-config} override and a new
\texttt{--output} destination. Cleanup restores the disabled environment
objects. Existing output directories are rejected, protecting the earlier
capture. The driver retains its explicit owned-server PID check and port-2000
guard. Local CARLA tooling and recorded sensor data remain untracked;
driver source, its patch, logs and native data are preserved as external
archive export copies. The configuration and this report are project artifacts.

\section{Executed validation}
An owned CARLA 0.10.0 instance on port 2011 ran Town10HD\_Opt with native
RGB, LiDAR, GNSS and IMU, using the same spawn index 105, Lincoln MKZ ego,
Dodge Charger target, 15 m initial gap and 11 m/s nominal waypoint-following
speeds as the previous capture. Fifty stationary calibration steps and
150 driving warm-up steps preceded 300 evaluation steps at 0.02 s.
The camera was 640-by-640 with a 100-degree field of view. LiDAR used
64 channels, 1,280,000 points/s, 50 revolutions/s, 80 m range and 0.02 m
range-noise standard deviation. LiDAR/IMU seed was 20261002 and GNSS seed
was 20261003. The default scene configuration was used without CLI override.

\begin{center}
\begin{tabular}{@{}lr@{}}
\toprule
Check & Result \\
\midrule
Disabled vehicle mesh components & 143 \\
Car / truck / bus / motorcycle / bicycle components & 47 / 30 / 2 / 8 / 56 \\
Native sample groups with aligned sensor frames & 500 / 500 \\
Saved evaluation RGB images / LiDAR clouds & 300 / 300 \\
LiDAR returns in the selected parked-car box, before / after & 140 / 0 \\
Cleanup errors / remaining vehicle and sensor actors & 0 / 0 \\
\bottomrule
\end{tabular}
\end{center}

The component count is not a count of distinct vehicles. At evaluation
frame 237, visual inspection confirmed removal of the roadside cars and
preservation of the red moving target. The ego and target positions exactly
matched the previous capture at this frame. Native LiDAR returns in the
oriented box of environment object 4120838475396380074 fell from 140 to zero.
Scoring-only transforms converted these raw points to world coordinates;
they were not perception or estimator inputs. All 500 simulation frame IDs
were consecutive. Capture took 75.327 s.

The first server launch exited with code 1 during startup without an explicit
cause in the retained log. A retry with a PTY and full stdout logging
succeeded; this does not establish the cause of the first exit. After capture,
a separate query confirmed zero vehicles and sensors and restored asynchronous
world settings. Only the owned server PID received SIGINT and exited with code
130. The shared port 2000 was not connected to, ticked, modified or signaled.

\section{Scope and reproducibility}
This validates scene configuration and native sensor capture. It does not
measure a new estimator RMSE or establish accuracy on scenes containing
parked traffic. Earlier recordings retain their original content. The
optional persistent association implementation remains available, but hiding
the parked objects is now a scene default independent of that option.

Compact results and exact executed commands are in
\path{report/CARLA_PARKED_VEHICLE_SCENE_20261002/}.
The local output is
\path{simulation_output/hide_parked_scene_20261002/native_capture/}.
\end{document}
